\section*{Hoofdstuk \ref{ch:g10:electron-shells} --- Elektronenschillen en het periodiek systeem}

\begin{solution}{exo:g10:electron-shells:1}
C: $1\mathrm{s}^2\,2\mathrm{s}^2\,2\mathrm{p}^2$. N:
$1\mathrm{s}^2\,2\mathrm{s}^2\,2\mathrm{p}^3$. Mg:
$1\mathrm{s}^2\,2\mathrm{s}^2\,2\mathrm{p}^6\,3\mathrm{s}^2$.
\end{solution}

\begin{solution}{exo:g10:electron-shells:2}
Koolstof 4, stikstof 5, magnesium 2.
\end{solution}

\begin{solution}{exo:g10:electron-shells:3}
Buitenste \omterm{def:g10:electron-shells:configuration}{schil} $n = 3$: periode 3. $2 + 3 = 5$ \omterm{def:g10:electron-shells:valence}{valentie-elektronen}:
\omterm{def:g9:periodic-table-first-look:period}{groep} 15 (fosfor).
\end{solution}

\begin{solution}{exo:g10:electron-shells:4}
\omterm{def:g10:electron-shells:family}{Alkalimetalen} (lithium, natrium, kalium); \omterm{def:g10:electron-shells:family}{halogenen} (fluor, chloor,
broom, jood); \omterm{def:g10:electron-shells:family}{edelgassen} (helium, neon, argon).
\end{solution}

\begin{solution}{exo:g10:electron-shells:5}
\omterm{def:g7:atoms-and-molecules:atom}{Atomen} hebben de neiging \omterm{def:g9:inside-the-atom:nucleus}{elektronen} op te nemen of af te staan zodat ze
acht \omterm{def:g9:inside-the-atom:nucleus}{elektronen} in hun buitenste \omterm{def:g10:electron-shells:configuration}{schil} krijgen, zoals het
dichtstbijzijnde \omterm{def:g10:electron-shells:family}{edelgas}.
\end{solution}

\begin{solution}{exo:g10:electron-shells:6}
K, $2.8.8.1$, staat er één af: \ce{K+},
$1\mathrm{s}^2\,2\mathrm{s}^2\,2\mathrm{p}^6\,3\mathrm{s}^2\,3\mathrm{p}^6$
(argon). F, $2.7$, neemt er één op: \ce{F-},
$1\mathrm{s}^2\,2\mathrm{s}^2\,2\mathrm{p}^6$ (neon).
\end{solution}

\begin{solution}{exo:g10:electron-shells:7}
$1\mathrm{s}^2\,2\mathrm{s}^2\,2\mathrm{p}^6\,3\mathrm{s}^2$: $Z = 12$,
magnesium.
\end{solution}

\begin{solution}{exo:g10:electron-shells:8}
\ce{Mg^{2+}} en \ce{Cl-}: \ce{MgCl2}. \ce{Al^{3+}} en \ce{O^{2-}}:
\ce{Al2O3}.
\end{solution}

\begin{solution}{exo:g10:electron-shells:9}
Zwavel (6 \omterm{def:g10:electron-shells:valence}{valentie-elektronen}, $2.8.6$), in dezelfde groep, 16.
\end{solution}

\begin{solution}{exo:g10:electron-shells:10}
Zijn buitenste \omterm{def:g10:electron-shells:configuration}{schil} is al vol (acht \omterm{def:g9:inside-the-atom:nucleus}{elektronen}): het heeft er niets bij
te winnen om \omterm{def:g9:inside-the-atom:nucleus}{elektronen} af te staan of op te nemen.
\end{solution}

\begin{solution}{exo:g10:electron-shells:11}
Argon heeft 18 \omterm{def:g9:inside-the-atom:nucleus}{elektronen}; \ce{X^{2+}} heeft er 2 afgestaan, dus X heeft
er 20: calcium.
\end{solution}

\begin{solution}{exo:g10:electron-shells:12}
Helium, $1\mathrm{s}^2$, heeft een volle buitenste \omterm{def:g10:electron-shells:configuration}{schil} (\omterm{def:g10:electron-shells:configuration}{schil} 1 bevat
maar twee \omterm{def:g9:inside-the-atom:nucleus}{elektronen}), net als de andere \omterm{def:g10:electron-shells:family}{edelgassen}; het gedraagt zich
niet als de \omterm{def:g9:periodic-table-first-look:metal}{metalen} van \omterm{def:g9:periodic-table-first-look:period}{groep} 2.
\end{solution}

\begin{solution}{exo:g10:electron-shells:13}
Zijn \omterm{def:g10:electron-shells:valence}{valentie-elektron} zit in \omterm{def:g10:electron-shells:configuration}{schil} 4, verder van de kern dan dat van
natrium in \omterm{def:g10:electron-shells:configuration}{schil} 3: het wordt minder sterk vastgehouden en gaat
makkelijker verloren.
\end{solution}

\begin{solution}{exo:g10:electron-shells:14}
Alle vier hebben 18 \omterm{def:g9:inside-the-atom:nucleus}{elektronen} en de configuratie van argon:
\[ 1\mathrm{s}^2\,2\mathrm{s}^2\,2\mathrm{p}^6\,3\mathrm{s}^2\,3\mathrm{p}^6. \]
\end{solution}

\begin{solution}{exo:g10:electron-shells:15}
Het neemt 3 \omterm{def:g9:inside-the-atom:nucleus}{elektronen} op om een octet te bereiken: het heeft 5
\omterm{def:g10:electron-shells:valence}{valentie-elektronen}, \omterm{def:g9:periodic-table-first-look:period}{groep} 15. In periode 3 is dat fosfor:
\[ 1\mathrm{s}^2\,2\mathrm{s}^2\,2\mathrm{p}^6\,3\mathrm{s}^2\,3\mathrm{p}^3. \]
\end{solution}

\begin{solution}{pb:g10:electron-shells:1}
\textbf{1.} Aluminium: 13 \omterm{def:g9:inside-the-atom:proton}{protonen}, 13 \omterm{def:g9:inside-the-atom:nucleus}{elektronen}. Zuurstof: 8 \omterm{def:g9:inside-the-atom:proton}{protonen},
8 \omterm{def:g9:inside-the-atom:nucleus}{elektronen}.

\textbf{2.} Al: $1\mathrm{s}^2\,2\mathrm{s}^2\,2\mathrm{p}^6\,3\mathrm{s}^2\,3\mathrm{p}^1$.
O: $1\mathrm{s}^2\,2\mathrm{s}^2\,2\mathrm{p}^4$.

\textbf{3.} Al: 3 \omterm{def:g10:electron-shells:valence}{valentie-elektronen}, periode 3, \omterm{def:g9:periodic-table-first-look:period}{groep} 13. O: 6,
periode 2, \omterm{def:g9:periodic-table-first-look:period}{groep} 16.

\textbf{4.} Aluminium is een \omterm{def:g9:periodic-table-first-look:metal}{metaal}, zuurstof een \omterm{def:g9:periodic-table-first-look:metal}{niet-metaal}.

\textbf{5.} \ce{Al^{3+}}: $1\mathrm{s}^2\,2\mathrm{s}^2\,2\mathrm{p}^6$.

\textbf{6.} \ce{O^{2-}}: $1\mathrm{s}^2\,2\mathrm{s}^2\,2\mathrm{p}^6$.

\textbf{7.} Neon.

\textbf{8.} Aluminium staat er 3 af, zuurstof neemt er 2 op.

\textbf{9.} $2 \times (+3) + 3 \times (-2) = 0$: \ce{Al2O3}.

\textbf{10.} Afgestaan: $2 \times 3 = 6$. Opgenomen: $3 \times 2 = 6$.

\textbf{11.} \omterm{def:g9:inside-the-atom:nucleus}{Elektronen} worden niet gemaakt of vernietigd: de \omterm{def:g9:inside-the-atom:nucleus}{elektronen}
die aluminium afstaat, zijn precies die welke zuurstof opneemt, en de
verbinding is neutraal.

\textbf{12.} Het heeft dezelfde lading, $+3$: de ladingen van het
kristal blijven in evenwicht.

\textbf{13.} $2 \times 27 + 3 \times 16 = 102$ \omterm{def:g9:inside-the-atom:proton}{nucleonen}.

\textbf{14.} $102 \times 1.67 \times 10^{-27} = \qty{1.70e-25}{kg}$.

\textbf{15.} $10^{-3} / 1.70 \times 10^{-25} \approx 5.87 \times 10^{21}$
formule-eenheden.

\textbf{16.} $2 \times 5.87 \times 10^{21} = 1.17 \times 10^{22}$
\ce{Al^{3+}}-ionen; $3 \times 5.87 \times 10^{21} = 1.76 \times 10^{22}$
\ce{O^{2-}}-ionen.

\textbf{17.} Een \omterm{def:g9:inside-the-atom:nucleus}{elektron} is ongeveer 1836 keer lichter dan een \omterm{def:g9:inside-the-atom:proton}{nucleon};
de 50 \omterm{def:g9:inside-the-atom:nucleus}{elektronen} van een formule-eenheid wegen minder dan drie
honderdste van één \omterm{def:g9:inside-the-atom:proton}{nucleon}.

\textbf{18.} $6 \times 5.87 \times 10^{21} \approx 3.5 \times 10^{22}$
\omterm{def:g9:inside-the-atom:nucleus}{elektronen} overgedragen voor \qty{1}{g} korund.
\end{solution}
